% Native author module 317. All three proof bodies retained in full.
% Original physical trajectories, events, labels, roof, conditioning and limits.
% Only scratch assembly; no inherited article or remote ref changed.

% Author proof development, 2026-10-19.  This file is not a referee report.
% Dependencies and the unclosed global-critical interface are explicit below.
\section{The actual endpoint charge in physical flux coordinates}
\label{sec:bf-actual-flux}

The physical radius, unit-speed clock, original section, exact collision,
return and lattice labels, conditioning by $c_*$, and incident-component
pairing are unchanged.  In this section $t=F$ is physical roof time,
$\theta=\arg\nabla_qF$ is its Gauss angle, and $u$ is a signed local
clearance parameter.  In particular $u$ is never substituted for $t$.
The original bad-source weight is $W^\varepsilon=1-H^\varepsilon$;
it equals one at a zero unvisited-disk clearance and is not set to zero
on that boundary.
For a positive clearance level $u_0>0$ we specify one incident side
of that level.  If both sides are counted, their charges must be
added explicitly.  The original unvisited tangency $u_0=0$ has the
prescribed physical side $u>0$.

Consider a one-sided regular physical marked-flight chart from
Theorem~\ref{thm:cf-flux}, with fixed flight index $j$, marked disk,
sign $\varsigma$, and actual collision word.  Put
\[
 q=\mathcal Q_j(\vartheta,u),\quad
 S(\vartheta,u)=F(q),\quad
 \theta(\vartheta,u)=\arg\nabla_qF(q).
\]
No reflection is made at the line's closest point.  At $u=0$ all
quantities below are one-sided limits of these actual positive-clearance
records.  In particular a limiting tangency is not counted as an extra
collision, and the collision index is not shifted.  Initially suppose
all visited collisions of the limiting word are nongrazing and the
endpoint gradient is nonzero.  The original endpoint cap is retained.

\begin{theorem}[Exact positive boundary flux]
\label{thm:bf-flux}
On a smooth boundary piece $u=u_0$ the original roof--angle trace is
\begin{equation}\label{eq:bf-flux}
 \frac{\omega W^\varepsilon}{4\pi c_*\chi}
       |\partial_\vartheta\theta|\,d\vartheta
   =\frac{W^\varepsilon}{4\pi R c_*}
       \left|\left.\frac{\partial u}{\partial t}\right|_\theta\right|
          d\vartheta .
\end{equation}
The formula is an identity of positive measures, with the actual
boundary multiplicity.  It therefore holds after a nonnegative Borel
restriction of already-defined incident boundary germs, including
their original record labels.  It does not assign a trace to a new
Borel selector that has no one-sided boundary value.

For an original section side at collision $i$, take
$(v,b)=(p_i,\alpha_i-a)$ or $(\alpha_i,p_i-a)$ on the corresponding
regular physical chart.  The same identity is
\begin{equation}\label{eq:bf-section-flux}
 \frac{\omega W^\varepsilon}{4\pi c_*\chi}
       |\partial_v\theta|\,dv
   =\frac{W^\varepsilon}{4\pi c_*}
       \left|\left.\frac{\partial b}{\partial t}\right|_\theta\right|
          dv .
\end{equation}
Thus the clearance formula has precisely one additional factor $R^{-1}$.
There is no additional conditional normalization in either formula.
\end{theorem}
\begin{proof}
Write $J=\det D_{(\vartheta,u)}(S,\theta)$ with its actual local sign.
The inverse of this two-by-two derivative matrix gives
\[
 \left.\frac{\partial u}{\partial t}\right|_\theta
       =-\frac{\partial_\vartheta\theta}{J}.
\]
The exact long-record determinant in \eqref{eq:cr-gauss} is
$|J|=\chi/(R\omega)$.  Substitute it into the preceding identity
and multiply by the physical line density $(4\pi Rc_*)^{-1}$.
This proves \eqref{eq:bf-flux}, including its radius factor.

The section formula follows in the same way from invariance of
$d\alpha\,dp$, rather than from a line-coordinate density.  Its
source density in $(v,b)$ is $(4\pi c_*)^{-1}$.  The ordinary
Jacobian chain rule gives $|\det D_{(v,b)}(t,\theta)|=\chi/\omega$.
Matrix inversion now gives \eqref{eq:bf-section-flux}.

The identities can first be tested on compact boundary subarcs where
the boundary-to-angle map is monotone.  Exhaust its open set of
nonzero derivative by countably many such subarcs, disjointizing their
source overlaps.  A $C^1$ real map sends its zero-derivative set to a
Lebesgue-null set: cover a compact portion of that set by intervals
on which the derivative has magnitude at most $\eta$, bound the
total image length by their total length times $\eta$, and let
$\eta\downarrow0$.  This also proves the needed one-dimensional
change of variables without an injectivity assumption on the full
boundary.  Nonnegative monotone convergence proves the identity for
the full indicated boundary family.

Distinct simultaneously active marked disks remain distinct marked
copies.  To represent a boundary counted only once by the original
pairing, choose one active constraint by a fixed finite ordering on
each finite physical word and restrict to its Borel incidence class.
The full marked sum is a positive upper bound before that restriction.
No physical collision word is multiplied by an unspecified chart count.
\end{proof}

\begin{proposition}[The transverse speed and its physical depth]
\label{prop:bf-speed}
Let $z=|\nabla_qF|$, $e=\nabla_qF/z$, and $G=D_q^2F$.
On a regular Gauss segment,
\begin{equation}\label{eq:bf-vector}
 V:=\left.\frac{dq}{dt}\right|_\theta
       =\frac{G^{-1}e}{z\,e^{\mathsf T}G^{-1}e},\qquad
 \|V\|_\infty\le\frac{\Lambda_\circ}{\lambda_\circ z}.
\end{equation}
For the clearance from any one of the actual nonincident disks of
flight $j$ whose foot is interior, at all regular points and one-sided
nongrazing limits,
\begin{equation}\label{eq:bf-speed}
 \left|\left.\frac{\partial u}{\partial t}\right|_\theta\right|
 \le \frac{2C_\Delta\Lambda_\circ}
               {\lambda_\circ R_-}
           \frac{\sigma^{d'_j}}{|(p_0,p_m)|},\qquad
 d'_j=\min(j,m-1-j).
\end{equation}
This is a bound on a physical normal speed, not an identity replacing
that speed by its upper bound near a critical point.
\end{proposition}
\begin{proof}
Differentiate $\nabla F=ze$ at fixed $e$, and use $dt/dz=z e^TG^{-1}e$.
The Hessian bounds in \eqref{eq:v109-hessian-bounds} give
\eqref{eq:bf-vector}.  The segment Hausdorff argument in
Proposition~\ref{prop:el-observables} applies separately to each fixed
candidate disk as well as to their minimum.  It gives
$|D_qu[\eta]|\le2C_\Delta\sigma^{d'_j}\|\eta\|_\infty$.
Combine this with \eqref{eq:bf-vector} and the exact identity
$z=R|(p_0,p_m)|$.  All differentiations are on the designated physical
word.  Continuity gives the stated one-sided nongrazing boundary limit.
\end{proof}

\subsection{A complete noncritical boundary estimate from the new
physical angular kernel}

The next result uses the all-time angular kernel of
Theorem~\ref{thm:nm-kernel} and its one-sided extension in
Theorem~\ref{thm:nm-trace}.  Their constant $H_{\rm nm}$ includes
both impact signs.  At $u=0$ this is the same one-sided unvisited
tangency family as above.  Both time directions are included.
This is an angular source measure, not the original conditioned
collision probability or an independent momentum law.

Let $\mathfrak B_{j,u_0}$ be the positive marked boundary charge in
\eqref{eq:bf-flux}, restricted to the original endpoint cap and any
original incident boundary labels, and put $r=|(p_0,p_m)|$.

\begin{theorem}[The full marked boundary away from the critical core]
\label{thm:bf-truncated}
For every
$0<\rho\le1$, every $u_0\in[0,3/200]$, and every finite physical
count,
\begin{equation}\label{eq:bf-truncated}
 \mathfrak B_{j,u_0}\{r\ge\rho\}
 \le C_{\rm bf}\,\sigma^{d'_j}
                   \bigl(1+\log(\sqrt2/\rho)\bigr),\qquad
 C_{\rm bf}=
 \frac{C_\Delta\Lambda_\circ H_{\rm nm}}
      {\pi c_*\lambda_\circ R_-^2}.
\end{equation}
Both signs and the complete true marked physical family are included.
Consequently for arbitrary choices $u_j\in[0,3/200]$,
\begin{equation}\label{eq:bf-depth-sum}
 \sum_{j:d'_j\ge D}\mathfrak B_{j,u_j}\{r\ge\rho\}
 \le\frac{2C_{\rm bf}\sigma^D}{1-\sigma}
                       \bigl(1+\log(\sqrt2/\rho)\bigr).
\end{equation}
In particular, for fixed $0<\vartheta_*<1/2$ and $a>0$, choosing
$D_m=\lceil\vartheta_*m\rceil$ and $\rho_m=e^{-am}$ in this
finite-count inequality gives convergence to zero after multiplication
by $m^2$.  This is convergence of a positive boundary charge in total
mass, and hence of its physical-roof pushforward in total variation.
It does not assert an $L^\infty$ roof density for that measure.
\end{theorem}
\begin{proof}
The complete two-sign kernel at the backward endpoint gives
\[
 \sum_{\varsigma=\pm1}\int\one_{\{r<s\}}\,d\vartheta
 \le\sum_{\varsigma=\pm1}\int\one_{\{|p_0|<s\}}\,d\vartheta
 \le2H_{\rm nm}s.
\]
The earlier endpoint is reached after $j+1\ge1$ actual collisions
backwards from the closest point.  Thus the kernel applies also to
$j=0$.  The complete angular measure satisfies
$A(s)\le2H_{\rm nm}s$.
Keeping original record tests only decreases this nonnegative measure.
With $L=\sqrt2$ and $r\le L$, the elementary identity
$r^{-1}=L^{-1}+\int_r^L s^{-2}\,ds$ gives
\[
 \int_{r\ge\rho}\frac{d\vartheta}{r}
 \le\frac{A(L)}L+\int_\rho^L\frac{A(s)}{s^2}\,ds
 \le2H_{\rm nm}\bigl(1+\log(L/\rho)\bigr).
\]
Use \eqref{eq:bf-flux}, $0\le W^\varepsilon\le1$, and
\eqref{eq:bf-speed}.  The prefactor is
$(4\pi R_-c_*)^{-1}(2C_\Delta\Lambda_\circ/
(\lambda_\circ R_-))\,2H_{\rm nm}$, which is the displayed
$C_{\rm bf}$.  This fixes both radius factors and the single $c_*^{-1}$.
No independence of the two endpoint momenta has been used.

There are at most two flight indices of each depth, so summing the
geometric series proves \eqref{eq:bf-depth-sum}.  At the stated finite
choices of $D_m,\rho_m$, the right side after multiplication by $m^2$
is bounded by a constant times
$m^2\sigma^{\vartheta_*m}(1+am+\log\sqrt2)$, tending to zero.
This uses an inequality valid for all finite $D,\rho,m$ and does not
change a microscopic roof window or the original order of outer and
collision limits.  The entire omitted critical core still remains an
original positive source; no claim about its smallness is made here.
\end{proof}

\subsection{Cancellation of the apparent pole at an actual critical
boundary germ}

\begin{theorem}[A finite angular trace at a nongrazing critical germ]
\label{thm:bf-critical-cancellation}
Let $q(s)$ be a $C^2$ parametrization of a one-sided smooth physical
boundary germ with $q(0)=q_*$, $q'(0)\ne0$, and all visited
incidences of the limiting word positive.  An unvisited tangency or an
original regular section side is allowed.  Assume $\nabla_qF(q_*)=0$.
Use the smooth extension of this designated stationary word only to
take one-sided derivatives; do not interpret its nonphysical side as
an actual trajectory.  Put
\[
 a=G_*q'(0),\qquad
 b=G_*q''(0)+D_q^3F(q_*)[q'(0),q'(0),\,\cdot\,].
\]
Then $a\ne0$, and the derivative of every continuous local lift of
the Gauss angle on either punctured side has the finite limit
\begin{equation}\label{eq:bf-critical-limit}
 \lim_{s\to0}\frac{d}{ds}\arg\nabla_qF(q(s))
       =\frac{\det(a,b)}{2|a|^2}.
\end{equation}
Consequently the positive boundary charge in
\eqref{eq:bf-flux} or \eqref{eq:bf-section-flux}, when expressed in
this boundary parameter, has a finite one-sided density at the germ.
The limiting density is
\begin{equation}\label{eq:bf-critical-density}
 \frac{\omega_*W_*^\varepsilon}
 {4\pi c_*\,e_a^{\perp\mathsf T}G_*e_a^\perp}
          \frac{|\det(a,b)|}{2|a|^2},\qquad e_a=a/|a|.
\end{equation}
The two antipodal limits of the Gauss angle have the same Hessian
denominator.  This is a trace of a boundary curve, not the different
incoming critical-fan coefficient of
Theorem~\ref{thm:mc-critical-sectors}.
\end{theorem}
\begin{proof}
Write $g(s)=\nabla_qF(q(s))$.  The stationary word is smooth in a
neighborhood of the limiting nongrazing contact configuration by its
positive interior Hessian.  The physical one-sided derivatives
therefore give
\[
 g(s)=as+\tfrac12bs^2+o(s^2),\qquad
 g'(s)=a+bs+o(s).
\]
Since $G_*$ is positive definite and $q'(0)\ne0$, $a\ne0$.
The angular derivative is exactly $\det(g,g')/|g|^2$ away from zero.
Its numerator is $\tfrac12\det(a,b)s^2+o(s^2)$ and its denominator
is $|a|^2s^2+o(s^2)$.  Division proves
\eqref{eq:bf-critical-limit}.  The finite one-sided optical and
guard limits have already been proved in Theorem~\ref{thm:v109-traces}.
Substitute them into the original coefficient $f^\varepsilon$ and
use $\chi=e^{\perp T}Ge^\perp$.  Replacing $e$ by $-e$ does not
change this quadratic form.  This proves
\eqref{eq:bf-critical-density}.

For the clearance formula one can also see the cancellation before
expanding: the boundary defining function $B$ is constant along
$q(s)$, hence $DB_*q'(0)=0$.  As $s\to0$, the leading direction of
$G^{-1}e$ is parallel to $q'(0)$.  Thus the leading numerator in
$DB\,G^{-1}e/z$ vanishes.  The expansion above computes the remaining
finite term, and does not merely delete its apparent $1/z$ bound.
\end{proof}

\begin{proposition}[A quantitative local turning bound]
\label{prop:bf-local-turning}
Suppose the preceding germ is parametrized by endpoint-angle
arclength on a one-sided interval, with $|q'|=1$ and $|q''|\le K$.
Put $L_G=2C_F+\Lambda_\circ K$.  On that same physical boundary
interval, for $0<|s|\le\lambda_\circ/L_G$,
\begin{equation}\label{eq:bf-local-turning}
 \left|\frac{d\theta}{ds}\right|
       \le\frac{2\Lambda_\circ L_G}{\lambda_\circ^2}.
\end{equation}
If $L_G=0$, the angular derivative is zero there.  The constants
$C_F,\lambda_\circ,\Lambda_\circ$ are the completed physical
constants of Sections~\ref{sec:el-localization} and
\ref{sec:v109-traces}; the boundary curvature allowance $K$ must be
proved on the actual boundary and is not asserted uniform here.
\end{proposition}
\begin{proof}
The physical third-variation bound gives
$\|dG/ds\|_{\rm op}\le2C_F$.  Consequently
$|g''(s)|\le2C_F+\Lambda_\circ K=L_G$ and
$|g'(s)|\le\Lambda_\circ$.  As
$|g'(0)|=|G_*q'(0)|\ge\lambda_\circ$,
\[
 |g(s)-s g'(0)|\le L_Gs^2/2,\qquad
 |g(s)|\ge\lambda_\circ|s|/2
\]
on the stated interval.  Use the exact identity
$g(s)=\int_0^s g'(v)\,dv$ to bound
\[
 |\det(g(s),g'(s))|
 \le \int_0^{|s|}\Lambda_\circ L_G(|s|-v)\,dv
 =\tfrac12\Lambda_\circ L_Gs^2.
\]
Division by the squared lower bound for $|g(s)|$ proves
\eqref{eq:bf-local-turning}.  The negative side has the same proof
with integration orientation reversed.  If $L_G=0$, $g'$ is constant
and the determinant is zero identically.
\end{proof}

\paragraph{The uniform continuation of this cancellation.}
The angular kernel controls $\int\one_{\{r<s\}}d\vartheta$ only
by $O(s)$, which explains the logarithm in
\eqref{eq:bf-truncated}.  The local cancellation just proved has a
uniform extension on the actual zero-offset critical-clearance
family in Theorem~\ref{thm:zf-flux} below.  That theorem proves a
common formal gradient domain, cancels the quadratic radial factors
before estimating, and sums the complete actual angular source
without a word-count constant.  The unrestricted nonzero-offset and
uncovered boundary pieces remain a separate geometric obligation;
bulk weak $L^{4/3}$ integrability does not supply it.  The incoming
critical fan at $g=0$ also retains its distinct original source.


% Complete finite-trajectory optical identity.  No pressure or new kernel input.
\section{Two actual endpoint momenta from a marked free-flight foot}
\label{sec:bf-foot-split}

Fix the original radius $R$, $m\ge1$, a true flight $0\le j<m$,
and one nonincident marked disk with signed impact
$\varsigma(R+u)$, $0<u\le3/200$.  All $m$ collisions and all physical
flights retain their original indices.  The closest point is a place
to parametrize the same line, with no reflection and no added
collision.  The limiting case $u=0$ below always means the one-sided
limit of this actual positive-clearance family on a fixed word with
positive limiting visited incidences.

The two true endpoint momenta are $p_0,p_m$.  Put
$r=|(p_0,p_m)|$, and denote by $a_-^{\rm foot}$ and
$a_+^{\rm foot}$ the first distances from the foot in the backward and
forward directions.  By Lemma~\ref{lem:xf-standoff},
\[
 a_\circ=\sqrt{1827}/200\le a_\pm^{\rm foot}\le\tau_+.
\]
There are $j+1$ actual collisions from the foot to the initial
collision in the reversed flow and $m-j$ actual collisions from the
foot to the terminal collision in the forward flow.  Both counts
are at least one, including when $m=1$.

For either direction propagate the unfolded transverse displacement
and angle field to the incoming side of its endpoint collision.  With
the positive physical matrices of Theorem~\ref{thm:xf-jacobian}, write
\begin{equation}\label{eq:bf-half-product}
 M_\pm=P(\ell_{n-1})K_{n-1}\cdots P(\ell_1)K_1
                         P(a_\pm^{\rm foot})
       =\begin{pmatrix}A_\pm&B_\pm\\C_\pm&D_\pm\end{pmatrix}.
\end{equation}
Here $n$ is the corresponding positive half-count.  For $n=1$ the
product is just $P(a_\pm^{\rm foot})$.  The endpoint collision is
not included as a kick in this incoming-side product.  Let
$\kappa_\pm=1/R$ and $c_\pm>0$ be its physical curvature and normal
incidence, and define
\begin{equation}\label{eq:bf-ab}
 A_\pm^{p}=\kappa_\pm B_\pm+c_\pm D_\pm,\qquad
 B_\pm^{p}=\kappa_\pm A_\pm+c_\pm C_\pm,\qquad
 h_\pm=B_\pm^{p}/A_\pm^{p}.
\end{equation}
The superscript $p$ distinguishes these positive scalar momentum
derivatives from entries of $M_\pm$.

\begin{theorem}[An exact two-sided physical momentum determinant]
\label{thm:bf-foot-determinant}
All scalars in \eqref{eq:bf-ab} are positive and
\begin{equation}\label{eq:bf-half-ratio}
 \frac1{\tau_++R_+}\le h_\pm\le\frac1{a_\circ}.
\end{equation}
On the original marked chart $(\vartheta,u)$ there are locally fixed
signs $s_0,s_m\in\{-1,1\}$ such that
\begin{equation}\label{eq:bf-signed-derivatives}
 \begin{array}{ll}
 \partial_\vartheta p_0=s_0 A_-^p,&
 \partial_u p_0=\varsigma s_0 B_-^p,\\[2pt]
 \partial_\vartheta p_m=s_m A_+^p,&
 \partial_u p_m=-\varsigma s_m B_+^p.
 \end{array}
\end{equation}
Consequently
\begin{equation}\label{eq:bf-positive-determinant}
 \left|\det D_{(\vartheta,u)}(p_0,p_m)\right|
       =A_-^p A_+^p(h_-+h_+)>0.
\end{equation}
The signs cannot cancel the positive sum.  No lower cutoff on any
positive visited incidence is required.
\end{theorem}
\begin{proof}
The reflection matrices were derived from the genuine reflection rule
in Theorem~\ref{thm:xf-jacobian}.  All their entries are nonnegative
and their determinants are one.  A positive initial free-flight
length makes $B_\pm>0$, while $A_\pm,D_\pm>0$ throughout.  Hence
both scalars in \eqref{eq:bf-ab} are positive.

At a terminal incoming collision write the incoming direction as
$v^-=-cn+pt$ and use arclength displacement $\delta q=t\delta s$.
For a physical incoming transverse field $(J,J')$ one has
$J=-c\delta s$.  Differentiating $p=v^-\cdot t$ gives
\[
 \delta p=-c\delta\vartheta^-+\kappa c\delta s
          =-(\kappa J+cJ').
\]
Unfolding a reflection reverses the signs of both transverse field
coordinates together.  Thus the actual endpoint momentum is a
locally fixed sign times the positive row $(\kappa,c)$ applied to
the propagated field.

The physical foot is $-\varsigma(R+u)v_\vartheta^\perp$.  Its
$\vartheta$ derivative is parallel to $v_\vartheta$, so the forward
transverse initial field is $(0,1)$.  Its $u$ derivative is
$-\varsigma v_\vartheta^\perp$, giving $(-\varsigma,0)$.
In the reversed direction $-v_\vartheta$, the angular derivative
is again one, while the same transverse displacement is now
$(\varsigma,0)$.  The momentum of a reversed collision is the
negative of the original momentum.  This changes the entire
backward endpoint row by one common sign and not its column ratio.
Applying the incoming-side products proves
\eqref{eq:bf-signed-derivatives}.  Taking their determinant proves
\eqref{eq:bf-positive-determinant}.

It remains to prove the uniform ratios without using a condition
number of the whole product.  Write a half-product as
$M=L P(a^{\rm foot})$.  If
$(X,Y)=(\kappa,c)L$, then $X>0$, $Y\ge0$, and
\[
 h=\frac{X}{a^{\rm foot}X+Y}
       =\frac1{a^{\rm foot}+Y/X}.
\]
For one collision, $L=I$, so $Y/X=c/\kappa\le R_+$.
For at least two collisions, the rightmost factor of $L$ is the
first actual reflection kick $K_1$.  Write
$(x,y)=(\kappa,c)\widehat L$, where $L=\widehat L K_1$.
Then $x>0$, $y\ge0$ and
\[
 (X,Y)=(x+y k_1,y),\qquad k_1=2/(Rc_1),\qquad
 \frac YX\le\frac1{k_1}\le\frac{R_+}{2}.
\]
This includes $y=0$, for which the ratio is zero.  The standoff
and horizon bounds now give \eqref{eq:bf-half-ratio}.  The proof
never bounds a grazing kick above; a smaller positive incidence
only increases that positive kick.
\end{proof}

\begin{corollary}[The exact endpoint kernel weight and its Wronskian]
\label{cor:bf-wronskian}
Let $u_x=|\partial_{\alpha_0}p_m|^{-1}$ be the actual endpoint
momentum density weight in Theorem~\ref{thm:rc-coarea}.  On the same
physical marked chart,
\begin{equation}\label{eq:bf-root-weight}
 u_x=\frac1{R A_-^p A_+^p(h_-+h_+)},\qquad
 u_x\le\frac{\tau_++R_+}{2R A_-^p A_+^p}.
\end{equation}
For $r>0$ and $\theta=\arg(-p_0,p_m)$ one has the exact identity
\begin{equation}\label{eq:bf-wronskian}
 u_x|\partial_\vartheta\theta|
 =\frac{
       \bigl|s_m p_0/A_-^p-s_0 p_m/A_+^p\bigr|}
        {R(h_-+h_+)r^2}.
\end{equation}
The original angle--roof trace is therefore
\begin{equation}\label{eq:bf-wronskian-trace}
 f^\varepsilon|\partial_\vartheta\theta|\,d\vartheta
 =\frac{W^\varepsilon}{4\pi c_*\,h_x}
   \frac{\bigl|s_m p_0/A_-^p-s_0 p_m/A_+^p\bigr|}
        {R(h_-+h_+)r^2}\,d\vartheta,
 \qquad h_x=\frac{R^2\chi}{\det G}.
\end{equation}
This uses the same $h_x$ and normalization as the exact radial
coarea theorem.  In particular $h_-+h_+$, a sum of the two half-ray
ratios, is not confused with $h_x$, the endpoint-momentum metric.
\end{corollary}
\begin{proof}
The one-copy line measure is
$d\nu=(4\pi R)^{-1}d\vartheta\,du$.  The actual endpoint-momentum
measure is $u_x(4\pi)^{-1}dp_0\,dp_m$, by
Theorem~\ref{thm:rc-coarea}.  Apply the local determinant
\eqref{eq:bf-positive-determinant} and compare the two actual
densities.  This proves the equality in \eqref{eq:bf-root-weight}.
Its inequality uses $h_-+h_+\ge2/(\tau_++R_+)$.

The derivative of the polar angle of $(-p_0,p_m)$ has absolute
value $|p_0\partial_\vartheta p_m-
p_m\partial_\vartheta p_0|/r^2$.  Substitute
\eqref{eq:bf-signed-derivatives} and cancel the product
$A_-^p A_+^p$ by \eqref{eq:bf-root-weight}.  This gives
\eqref{eq:bf-wronskian}.  Finally the exact determinant of
$q\mapsto(p_0,p_m)$ is $\det G/R^2$, so
$u_x=R^2\omega/\det G$.  Hence
$f^\varepsilon=W^\varepsilon u_x/(4\pi c_*h_x)$, proving
\eqref{eq:bf-wronskian-trace}.  All these comparisons were between
the same unconditioned collision measure before the one factor
$c_*^{-1}$ was inserted.
\end{proof}

\paragraph{What the identity supplies.}
Equations~\eqref{eq:bf-root-weight}--\eqref{eq:bf-wronskian-trace}
are unconditional identities and bounds on the actual regular
physical marked family; they do not use a whole-word pressure
theorem, mixing, or a near-miss kernel bound.  They keep the
potentially useful signed cancellation between the two endpoint
terms.  Bounding those two terms separately recovers a possible
$1/r$ loss and does not prove the complete critical boundary sum.
At a shared actual critical germ, both signed first-order terms in
the numerator agree, and their difference cancels.  A uniform
weighted summation of that cancellation over arbitrary short
physical boundary components is a further estimate, not a
consequence of the positive determinant alone.


% Complete author proof.  The formal domain is used only for analytic identities;
% the integrated source consists exclusively of actual one-sided physical records.
\section{A uniform cancellation on actual critical-clearance boundary pieces}
\label{sec:zf-zero-offset}

The physical model and all exact record selectors remain unchanged.
The result below concerns a complete explicitly specified part of the
original one-sided clearance boundary.  It uses a common domain of
the formal reflection equations already proved in
Theorem~\ref{thm:pt-analytic}; it does not declare that formal domain
to be a physical neighborhood or to preserve a section test.
For $u_0>0$ fix one incident side of the level; a two-sided charge
is the explicit sum of two such measures.  At $u_0=0$ the physical
side is the original positive-clearance side.

Keep $\lambda=50/47$, and put $\gamma=\lambda^{-1}=47/50$.
Let $\eta=\eta_*$ and $P=P_*$ be exactly the analytic constants in
\eqref{eq:pt-analytic-constants}.  Define the following positive
constants, all independent of collision count and word:
\begin{align}
 \underline g&=R_-/2=9/40,&
 L_z&=28/5+R_++3/200=1217/200,\notag\\
 K_F^{\rm an}&=2048R_+/\eta^2,&
 a_{\rm an}&=\min\{\eta/16,\underline g/(2K_F^{\rm an})\},
                                                     \label{eq:zf-constants}\\
 b_{\rm an}&=\underline g\,a_{\rm an}/4,&
 L_B&=P(4R_++12(R_+/\tau_-)L_z),\notag\\
 K_{B,1}&=6L_B,&K_{B,2}&=128L_B/\eta,\notag\\
 K_{B,g}&=K_{B,2}/\underline g^2
             +K_{B,1}K_F^{\rm an}/\underline g^3,&
 C_{\rm zf}&=\Lambda_\circ K_{B,g}/(R_-c_*).
                                                     \label{eq:zf-final-constant}
\end{align}
The constant $28/5$ is the enclosing true-flight bound in
Theorem~\ref{thm:eh-horizon}; it does not replace the intrinsic
physical horizon or the physical clock.

Fix a limiting physical critical configuration $Q_*$ of a designated
$m$-flight word, with endpoint gradient zero and the original endpoint
cap.  It is the limit of actual regular trajectories of that word.
Its endpoint incidences are one, and any number of interior incidences
may be zero.  Fix a nonincident marked disk on flight $j$ whose
perpendicular foot is interior and whose limiting signed impact is
$\varsigma(R+u_0)$, $u_0\in[0,3/200]$.
Let $q_*=(\alpha_0^*,\alpha_m^*)$ in fixed local lifts.
The formal analytic reflection lift based at this configuration
defines the formal action $F$ and the signed line-clearance function
\begin{equation}\label{eq:zf-B}
 B(q)=\varsigma\,(z_{\rm mark}-x_j(q))\cdot v_j(q)^\perp-R-u_0.
\end{equation}
Thus $\nabla_qF(q_*)=0$ and $B(q_*)=0$ exactly.
At every actual selected point the same expression is the true
clearance from that marked disk minus $u_0$.  Its use at intermediate
formal points is an analytic calculation, not a visibility assertion.

\begin{lemma}[Uniform formal gradient disk and signed-clearance tensors]
\label{lem:zf-formal}
On the real endpoint ball $|q-q_*|<a_{\rm an}$ the formal Hessian
satisfies
\begin{equation}\label{eq:zf-hessian}
 \underline g I\preceq G(q)\preceq2\Lambda_\circ I,
 \qquad \|D_qG(q)[h]\|_{\rm op}
                           \le K_F^{\rm an}|h|.
\end{equation}
Every $|g|<b_{\rm an}$ has a unique formal inverse $q(g)$ in this
ball, and it lies in $|q(g)-q_*|<a_{\rm an}/4$.  Put
$\mathcal B(g)=B(q(g))$.  For all these gradients,
\begin{equation}\label{eq:zf-Bg}
 \|D_g^2\mathcal B(g)\|_{\rm op}
             \le K_{B,g}\,\gamma^{d'_j},\qquad
 d'_j=\min(j,m-1-j).
\end{equation}
No positive lower bound on an interior incidence or on an unvisited
clearance is required for this formal statement.
\end{lemma}
\begin{proof}
The limiting-reference version of Theorem~\ref{thm:pt-analytic}
gives a holomorphic lift on $|h_0|+|h_m|<\eta$.  Consider the
polydisk $|h_0|,|h_m|<\eta/4$.  Its weighted contact increment is
at most $P\eta/2=r_*/8$.  All flight complexifications remain in
the nonvanishing disks of Lemma~\ref{lem:pt-complex}.  More explicitly,
the two complex components $X^+,X^-$ differ from their real reference
values by at most a quarter of the reference flight length.  Choose
the square root of $X^+X^-$ that is positive at the reference point.
Its magnitude is at least three quarters of that reference length,
and the holomorphic unit direction has complex Euclidean norm at
most $5/3$.  Circular normals and tangents have norm less than two
on the contact ball.  The endpoint first-variation identity therefore
gives a holomorphic gradient with each component bounded by $4R_+$.
This bound is independent of $m$: it is a bound for
$R(-p_0,p_m)$, not for the total action.

For completeness, the first-variation identity continues to the
formal lift because its reflection equations imply vanishing interior
tangential derivatives of the same analytic action.  It can also be
obtained by holomorphic continuation of the identity on the original
regular branch.  No intermediate formal contact is asserted physical.

Cauchy's formula on the smaller polydisk
$|h_0|,|h_m|<\eta/8$ bounds every second partial derivative of each
gradient component by
$2(4R_+)(8/\eta)^2=512R_+/\eta^2$.
For a two-by-two Hessian derivative, summing its two entries in each
row and the two input coordinates bounds its Euclidean operator norm
by four times this number times $|h|$.  This is precisely
$K_F^{\rm an}$ in \eqref{eq:zf-constants}.

At the critical physical limit $G_*\succeq R I\succeq
2\underline g I$, by the continuous reduced Hessian formula, and
$G_*\preceq\Lambda_\circ I$.  Integrating the preceding derivative
bound along a real endpoint segment of length at most $a_{\rm an}$
changes the Hessian by at most $\underline g/2$.  This proves
\eqref{eq:zf-hessian}, with room in both inequalities.
For a prescribed $|g|<b_{\rm an}$, minimize
$F(q_*+h)-g\cdot h$ on $|h|\le a_{\rm an}$.
Strict convexity holds by \eqref{eq:zf-hessian}.  On the boundary its
outward radial derivative is at least
$\underline g a_{\rm an}^2-|g|a_{\rm an}>0$.
The minimizer is therefore interior and is the unique inverse of the
gradient in that ball.  Strong monotonicity then gives
$|q(g)-q_*|\le |g|/\underline g<a_{\rm an}/4$.
The inverse function theorem makes it smooth on the whole gradient
disk, and differentiation yields
\begin{equation}\label{eq:zf-inverse-tensors}
 \|D_gq\|_{\rm op}\le\underline g^{-1},\qquad
 \|D_g^2q\|_{\rm op}\le K_F^{\rm an}\underline g^{-3}.
\end{equation}

We next check the depth factor for the individual marked clearance.
This is a bound for its signed line coordinate, so no differentiability
of a least-clearance minimizer is assumed.  Let $s=|h_0|+|h_m|$ and
$d=d'_j$.  The weighted analytic localization gives
\[
 |\alpha_j(h)-\alpha_j^*|
       +|\alpha_{j+1}(h)-\alpha_{j+1}^*|
                       \le P s\gamma^d.
\]
The corresponding complex position difference is at most
$2R_+Ps\gamma^d$.  Along the straight contact increment, all
complex direction differentials have norm at most
\[
 \frac{1+(5/3)^2}{3\tau_-/4}\,|dX|
       =\frac{136}{27\tau_-}|dX|
       <\frac6{\tau_-}|dX|.
\]
Thus the flight-direction difference is at most
$12(R_+/\tau_-)Ps\gamma^d$.  The reference vector from departure
contact to the marked center has length at most $L_z$.  In
\eqref{eq:zf-B}, bound the position increment against the new
direction, whose norm is less than two, and the direction increment
against this fixed reference vector.  The result is
\[
 |B(q_*+h)-B(q_*)|\le L_Bs\gamma^d
                   \le L_B\eta\gamma^d/2
\]
on the larger polydisk.  Cauchy's formula on the smaller polydisk
therefore gives Euclidean derivative bounds
\[
 |D_qB|\le6L_B\gamma^d=K_{B,1}\gamma^d,
 \qquad
 \|D_q^2B\|_{\rm op}
          \le128L_B\eta^{-1}\gamma^d=K_{B,2}\gamma^d.
\]
The first constant bounds $4\sqrt2L_B$; for the second, each
second partial is at most $64L_B/\eta$, and the two-by-two operator
norm is at most twice that bound.  Composing these inequalities with
\eqref{eq:zf-inverse-tensors}, keeping both second-order chain-rule
terms, proves exactly \eqref{eq:zf-Bg}.
\end{proof}

Call an actual one-sided boundary point \emph{covered by a critical
zero-offset chart} if its designated word and marked disk agree with
one of the preceding limiting critical configurations, it has
$B(q)=0$, and its full contact configuration agrees with that formal
inverse $q(g)$ at some $0<|g|<b_{\rm an}$.  Agreement here includes
all $m+1$ contacts, not only the two endpoint angles.  The point must
itself be an actual regular one-sided boundary record with all
original physical cuts retained.  The formal inverse supplies no new
collision event.  This definition contains the actual portions of
every such critical boundary germ within the common formal gradient
disk, irrespective of their possibly short or disconnected physical
continuation.  It does not include a branch whose corresponding
formal critical clearance is nonzero.

\begin{theorem}[Uniform full angular-source flux on the covered boundary]
\label{thm:zf-flux}
For every actual covered boundary point,
\begin{equation}\label{eq:zf-speed}
 \left|\left.\frac{\partial u}{\partial t}\right|_\theta\right|
                    \le\Lambda_\circ K_{B,g}\gamma^{d'_j}.
\end{equation}
Consequently the complete original positive marked endpoint charge
on all covered boundary points at flight $j$, with both signs and
any retained original exact labels, satisfies
\begin{equation}\label{eq:zf-flux}
 \mathfrak B^{\rm crit}_{j,u_0}
                    \le C_{\rm zf}\gamma^{d'_j}.
\end{equation}
There is no deletion of small positive $|(p_0,p_m)|$ in this sum,
no count of physical words in its constant, and no angular-kernel or
whole-word pressure assumption.  Summing all flights of depth at
least $D$ gives
\begin{equation}\label{eq:zf-deep}
 \sum_{j:d'_j\ge D}\mathfrak B^{\rm crit}_{j,u_j}
           \le\frac{2C_{\rm zf}}{1-\gamma}\gamma^D,
 \qquad u_j\in[0,3/200].
\end{equation}
For every fixed $0<\vartheta_*<1/2$, its product with $m^2$ tends
to zero at $D=\lceil\vartheta_*m\rceil$, uniformly in the original
radius and labels.  These are bounds for the positive endpoint
charge and for the total variation of its pushforward by the
unchanged physical roof, not an assertion of a roof density for
that boundary measure.
The incoming critical fan at $g=0$ itself is not this boundary-curve
charge.  Its original angular coefficient remains the separate
quantity in Theorem~\ref{thm:mc-critical-sectors}; it is not assigned
zero or included in \eqref{eq:zf-flux} by a bulk-null convention.
\end{theorem}
\begin{proof}
Write $g=\nabla_qF$ at the selected actual point.  Both
$\mathcal B(0)=0$ and $\mathcal B(g)=0$ hold.  The entire straight
segment from $0$ to $g$ lies in the \emph{formal} gradient disk of
Lemma~\ref{lem:zf-formal}, where the proved analytic derivative
bound is valid.  The exact integral remainder is
\begin{equation}\label{eq:zf-euler}
 D_g\mathcal B(g)\cdot g
 =\mathcal B(g)-\mathcal B(0)
       +\int_0^1 sD_g^2\mathcal B(sg)[g,g]\,ds
 =\int_0^1 sD_g^2\mathcal B(sg)[g,g]\,ds.
\end{equation}
Its absolute value is at most
$K_{B,g}\gamma^{d'_j}|g|^2/2$.
At the actual point, the roof vector at fixed Gauss angle is
$G^{-1}g/(g^{\mathsf T}G^{-1}g)$, so the normal speed is precisely
\[
 \left.\frac{\partial u}{\partial t}\right|_\theta
  =\frac{D_g\mathcal B(g)\cdot g}
            {g^{\mathsf T}G^{-1}g}.
\]
Since the denominator is at least $|g|^2/(2\Lambda_\circ)$,
the quadratic factors cancel and prove \eqref{eq:zf-speed}.
This calculation evaluates the physical normal speed at an actual
point; it makes no physical-flow claim about the formal segment
used in the integral identity.

The exact marked boundary formula \eqref{eq:bf-flux} now bounds
the charge by
\[
 \frac{\Lambda_\circ K_{B,g}\gamma^{d'_j}}
      {4\pi R_-c_*}
        \sum_{\varsigma=\pm1}\int_0^{2\pi}d\vartheta
 = C_{\rm zf}\gamma^{d'_j}.
\]
For each angular parameter and sign the limiting original finite
record supplies only one marked physical source.  Its possible
assignment to a critical chart is made on this source before taking
the upper bound.  The bound is pointwise and identical for every
such chart; it is not multiplied by their number.

Here is a precise disjoint assignment if multiple coordinate lifts
were retained.  At fixed $R,m$ there are finitely many physical words
and finitely many candidate marked disks per flight.  A given word
has at most one limiting normal-endpoint physical configuration.
Indeed the support inequalities of Lemma~\ref{lem:v111-support}
pass to its one-sided limits.  If both endpoint gradients are zero,
the endpoint normal incidences are one, so the summed inequality
forces both endpoint contacts to agree.  Both convex support gaps
then vanish.  Equality in the norm support inequality makes every
pair of corresponding flight vectors positively collinear.  Starting
at the common first contact, the incoming-root condition
$c_i\ge0$ chooses the same entry or tangent point of each next
designated disk, successively.  All contacts therefore agree,
including possible limiting visited tangencies.  Thus only finitely
many reference configurations occur.  Each of their formal inverse
charts is continuous, and its equality with the actual finite record
and all original inequalities is Borel.  Assign an overlapping
source point to the first such chart in a fixed ordering.  Artificial
angular seams are assigned once.  The nonnegative source bound is
unchanged.  This proves the complete sum rather than a separate
estimate per word.

There are at most two flight indices at each depth, and
$\sum_{d\ge D}\gamma^d=\gamma^D/(1-\gamma)$.
This proves \eqref{eq:zf-deep}.  Multiplication by $m^2$ at the
stated finite choice of $D$ leaves a polynomial times an exponential
with strictly negative exponent.  No microscopic roof interval or
outer limit has been changed.
\end{proof}

\paragraph{The exact remaining nonzero-offset term.}
This theorem treats a geometric class by a proved uniform normal-speed
cancellation, not by assigning its positive trace zero.  At an actual
boundary point represented in a formal gradient chart with
$\mathcal B(0)=a_w\ne0$, the same exact calculation instead gives
\[
 \left.\frac{\partial u}{\partial t}\right|_\theta
 =\frac{-a_w+\int_0^1sD_g^2\mathcal B(sg)[g,g]\,ds}
        {g^{\mathsf T}G^{-1}g}.
\]
The offset is an actual signed clearance of that specified formal
critical word, not a freely chosen clock or threshold.  It cannot be
dropped even if small.  A complete weighted angular-source estimate
for this offset term, for uncovered critical pieces, and for the
original section/pairing boundary terms remains required to close
the entire endpoint charge.  The theorem neither interchanges a
local critical-radius limit with $m\to\infty$ nor concludes the
original essential-height or raw pointwise local limit assertion.
